\documentclass[a4paper,12pt]{article}

\usepackage{polyglossia} 
\setmainlanguage{russian}
\usepackage{fontspec}
\setmainfont{Times New Roman}
\usepackage{geometry}
\usepackage{amsmath}
\usepackage{graphicx} 
\usepackage{listings}
\usepackage{xcolor}

\geometry{top=20mm, bottom=20mm, left=25mm, right=15mm}

\lstset{
    language=Python,
    basicstyle=\ttfamily\scriptsize,
    keywordstyle=\color{blue},
    stringstyle=\color{red},
    breaklines=true,
    frame=single
}

\begin{document}

\noindent
\begin{minipage}[c]{0.22\textwidth}
    \centering
    
\includegraphics[width=\textwidth]{../bmstu_logo.jpg} 
\end{minipage}%
\hfill
\begin{minipage}[c]{0.75\textwidth}
    \centering
    {\small\bfseries Министерство науки и высшего образования Российской Федерации} \\[2pt]
    {\small\bfseries Федеральное государственное автономное образовательное учреждение} \\[2pt]
    {\small\bfseries высшего образования} \\[2pt]
    {\small\bfseries «Московский государственный технический университет} \\[2pt]
    {\small\bfseries имени Н.Э. Баумана} \\[2pt]
    {\small\bfseries (МГТУ им. Н.Э. Баумана)}
\end{minipage}

\vspace{12pt}
\hrule height 1pt
\vspace{2pt}
\hrule height 2pt

\vspace{10mm}

\noindent
\textbf{ФАКУЛЬТЕТ} \underline{\makebox[138mm][l]{\hspace{15mm}\textbf{ИНФОРМАТИКА И СИСТЕМЫ УПРАВЛЕНИЯ}}} \\[10pt]
\textbf{КАФЕДРА} \underline{\makebox[143mm][l]{\hspace{12mm}\textbf{СИСТЕМЫ ОБРАБОТКИ ИНФОРМАЦИИ И УПРАВЛЕНИЯ}}}

\vspace{25mm}

\begin{center}
    {\Large Отчет по лабораторной работе № 9} \\[12pt]
    {\Large\bfseries «Применение глубокого обучения.} \\[4pt]
    {\Large\bfseries Архитектуры нейронных сетей»} \\[20mm]
    {\large по дисциплине «Технологии машинного обучения»}
\end{center}

\vspace{35mm}

\noindent
\begin{tabular}{@{}llcc}
    \large Студент & \underline{\makebox[28mm][c]{\textbf{ИУ5-64Б}}} & \underline{\makebox[55mm]{}} & \underline{\makebox[50mm][c]{\textbf{Н.А. Чернев}}} \\[2pt]
    & & & \scriptsize (И.О. Фамилия) \\[20pt]
    \large Преподаватель & & \underline{\makebox[55mm]{}} & \underline{\makebox[50mm][c]{\textbf{Ю.Е. Гапанюк}}} \\[2pt]
    & & & \scriptsize (И.О. Фамилия)
\end{tabular}

\vfill
\begin{center}
    \large 2026 г.
\end{center}

\newpage

\section*{Цель работы}
Изучение архитектур глубоких нейронных сетей и их применения к задачам классификации/регрессии.

\section*{Подготовка данных}

\begin{lstlisting}
import numpy as np
import tensorflow as tf
from tensorflow import keras
from keras import layers
from sklearn.datasets import load_iris
from sklearn.preprocessing import StandardScaler
from sklearn.model_selection import train_test_split

# Загрузка данных
X, y = load_iris(return_X_y=True)

# Нормализация
scaler = StandardScaler()
X_scaled = scaler.fit_transform(X)

# Разделение
X_train, X_test, y_train, y_test = train_test_split(
    X_scaled, y, test_size=0.3, random_state=42)

# One-hot кодирование целевой переменной
y_train_cat = keras.utils.to_categorical(y_train, 3)
y_test_cat = keras.utils.to_categorical(y_test, 3)
\end{lstlisting}
\section*{Подготовка данных}

\begin{lstlisting}
import numpy as np
import tensorflow as tf
from tensorflow import keras
from keras import layers
from sklearn.datasets import load_iris
from sklearn.preprocessing import StandardScaler
from sklearn.model_selection import train_test_split
from sklearn.decomposition import PCA
import matplotlib.pyplot as plt

# Загрузка данных
X, y = load_iris(return_X_y=True)

# Нормализация
scaler = StandardScaler()
X_scaled = scaler.fit_transform(X)

# Разделение
X_train, X_test, y_train, y_test = train_test_split(
    X_scaled, y, test_size=0.3, random_state=42)

# One-hot кодирование целевой переменной
y_train_cat = keras.utils.to_categorical(y_train, 3)
y_test_cat = keras.utils.to_categorical(y_test, 3)
\end{lstlisting}

\subsection*{Анализ главных компонент для визуализации}

\begin{lstlisting}
# PCA для 2D визуализации
pca = PCA(n_components=2)
X_pca = pca.fit_transform(X_scaled)

# Вклад компонент
fig, ax = plt.subplots(figsize=(8, 6))
components = ['PC1', 'PC2']
variance = pca.explained_variance_ratio_ * 100
ax.bar(components, variance, color=['#FF6B6B', '#4ECDC4'])
ax.set_ylabel('Объясненная дисперсия (%)')
ax.set_title('Вклад главных компонент')
plt.savefig('fig_1.png', dpi=150, bbox_inches='tight')
plt.show()
\end{lstlisting}

\begin{figure}[h!]
    \centering
    \includegraphics[width=0.85\textwidth]{fig_1.png}
    \small Рисунок 1 --- Вклад главных компонент в объясненную дисперсию
\end{figure}

\subsection*{Визуализация датасета в 2D пространстве}

\begin{lstlisting}
# t-SNE редукция для лучшей визуализации
from sklearn.manifold import TSNE

tsne = TSNE(n_components=2, perplexity=30, n_iter=1000, random_state=42)
X_tsne = tsne.fit_transform(X_scaled)

# Визуализация датасета после t-SNE
fig, ax = plt.subplots(figsize=(10, 8))
colors = ['#FF6B6B', '#4ECDC4', '#45B7D1']
for i in range(3):
    mask = y == i
    ax.scatter(X_tsne[mask, 0], X_tsne[mask, 1], 
               c=colors[i], label=f'Class {i}', s=80, alpha=0.7)
ax.set_xlabel('t-SNE компонента 1')
ax.set_ylabel('t-SNE компонента 2')
ax.set_title('t-SNE редукция датасета Iris')
ax.legend()
plt.savefig('fig_2.png', dpi=150, bbox_inches='tight')
plt.show()
\end{lstlisting}

\begin{figure}[h!]
    \centering
    \includegraphics[width=0.85\textwidth]{fig_2.png}
    \small Рисунок 2 --- t-SNE редукция датасета (локальная структура)
\end{figure}

\section*{Архитектуры нейронных сетей}

\subsection*{Мелкая сеть (2 слоя)}
\begin{lstlisting}
model_shallow = keras.Sequential([
    layers.Dense(128, activation='relu', input_shape=(4,)),
    layers.Dropout(0.2),
    layers.Dense(64, activation='relu'),
    layers.Dropout(0.2),
    layers.Dense(3, activation='softmax')
])

model_shallow.compile(optimizer='adam', 
                      loss='categorical_crossentropy',
                      metrics=['accuracy'])

history_shallow = model_shallow.fit(
    X_train, y_train_cat, epochs=100, batch_size=16, 
    validation_split=0.2, verbose=0)
\end{lstlisting}

\subsection*{Средняя сеть (3 слоя)}
\begin{lstlisting}
model_medium = keras.Sequential([
    layers.Dense(256, activation='relu', input_shape=(4,)),
    layers.Dropout(0.3),
    layers.Dense(128, activation='relu'),
    layers.Dropout(0.3),
    layers.Dense(64, activation='relu'),
    layers.Dropout(0.2),
    layers.Dense(3, activation='softmax')
])

model_medium.compile(optimizer='adam',
                     loss='categorical_crossentropy',
                     metrics=['accuracy'])

history_medium = model_medium.fit(
    X_train, y_train_cat, epochs=100, batch_size=16,
    validation_split=0.2, verbose=0)
\end{lstlisting}

\subsection*{Глубокая сеть (5 слоев)}
\begin{lstlisting}
model_deep = keras.Sequential([
    layers.Dense(512, activation='relu', input_shape=(4,)),
    layers.Dropout(0.4),
    layers.Dense(256, activation='relu'),
    layers.Dropout(0.3),
    layers.Dense(128, activation='relu'),
    layers.Dropout(0.3),
    layers.Dense(64, activation='relu'),
    layers.Dropout(0.2),
    layers.Dense(32, activation='relu'),
    layers.Dropout(0.2),
    layers.Dense(3, activation='softmax')
])

model_deep.compile(optimizer='adam',
                   loss='categorical_crossentropy',
                   metrics=['accuracy'])

history_deep = model_deep.fit(
    X_train, y_train_cat, epochs=100, batch_size=16,
    validation_split=0.2, verbose=0)
\end{lstlisting}

\section*{Анализ процесса обучения}

\subsection*{История обучения моделей}

\begin{lstlisting}
# Визуализация процесса обучения
fig, axes = plt.subplots(1, 2, figsize=(14, 5))

# Loss
for history, name in [(history_shallow, 'Shallow'),
                       (history_medium, 'Medium'),
                       (history_deep, 'Deep')]:
    axes[0].plot(history.history['loss'], label=name, linewidth=2)

axes[0].set_xlabel('Эпоха')
axes[0].set_ylabel('Loss')
axes[0].set_title('История обучения (Loss)')
axes[0].legend()
axes[0].grid(True, alpha=0.3)

# Accuracy
for history, name in [(history_shallow, 'Shallow'),
                       (history_medium, 'Medium'),
                       (history_deep, 'Deep')]:
    axes[1].plot(history.history['accuracy'], label=name, linewidth=2)

axes[1].set_xlabel('Эпоха')
axes[1].set_ylabel('Accuracy')
axes[1].set_title('История обучения (Accuracy)')
axes[1].legend()
axes[1].grid(True, alpha=0.3)

plt.tight_layout()
plt.savefig('fig_3.png', dpi=150, bbox_inches='tight')
plt.show()
\end{lstlisting}

\begin{figure}[h!]
    \centering
    \includegraphics[width=0.85\textwidth]{fig_3.png}
    \small Рисунок 3 --- История обучения (Loss и Accuracy) для разных архитектур
\end{figure}

\subsection*{Сравнение архитектур}

\begin{lstlisting}
# Тестирование моделей
models = {
    'Shallow (2 layers)': model_shallow,
    'Medium (3 layers)': model_medium,
    'Deep (5 layers)': model_deep
}

results = {}
for name, model in models.items():
    test_loss, test_acc = model.evaluate(X_test, y_test_cat, verbose=0)
    results[name] = test_acc
    print(f"{name}: Accuracy = {test_acc:.4f}")

# Визуализация результатов
fig, ax = plt.subplots(figsize=(10, 6))
names = list(results.keys())
accuracies = list(results.values())
colors = ['#FF6B6B', '#4ECDC4', '#45B7D1']
ax.bar(names, accuracies, color=colors)
ax.set_ylabel('Test Accuracy')
ax.set_title('Сравнение точности архитектур на тестовой выборке')
ax.set_ylim([0.9, 1.0])
for i, (name, acc) in enumerate(zip(names, accuracies)):
    ax.text(i, acc + 0.005, f'{acc:.4f}', ha='center')
plt.xticks(rotation=15, ha='right')
plt.tight_layout()
plt.savefig('fig_4.png', dpi=150, bbox_inches='tight')
plt.show()
\end{lstlisting}

\begin{figure}[h!]
    \centering
    \includegraphics[width=0.85\textwidth]{fig_4.png}
    \small Рисунок 4 --- Сравнение точности различных архитектур нейронных сетей
\end{figure}

\section*{Анализ предсказаний}

\subsection*{Предсказания разных моделей}

\begin{lstlisting}
# Получение предсказаний
y_pred_shallow = model_shallow.predict(X_test)
y_pred_medium = model_medium.predict(X_test)
y_pred_deep = model_deep.predict(X_test)

# Преобразование в классы
y_pred_shallow_cls = np.argmax(y_pred_shallow, axis=1)
y_pred_medium_cls = np.argmax(y_pred_medium, axis=1)
y_pred_deep_cls = np.argmax(y_pred_deep, axis=1)

# Визуализация результатов классификации для мелкой сети
fig, ax = plt.subplots(figsize=(10, 8))
for i in range(3):
    mask = y_test == i
    ax.scatter(X_pca[mask, 0], X_pca[mask, 1], 
               c='gray', s=100, alpha=0.3, marker='o', label=f'True Class {i}')
    
for i in range(3):
    mask = (y_test_cls == i) & (y_pred_shallow_cls == i)
    ax.scatter(X_pca[mask, 0], X_pca[mask, 1], 
               c=['#FF6B6B', '#4ECDC4', '#45B7D1'][i], s=50, marker='o')

ax.set_xlabel('PC1')
ax.set_ylabel('PC2')
ax.set_title('Предсказания Shallow Network в PCA пространстве')
plt.tight_layout()
plt.savefig('fig_5.png', dpi=150, bbox_inches='tight')
plt.show()
\end{lstlisting}

\begin{figure}[h!]
    \centering
    \includegraphics[width=0.85\textwidth]{fig_5.png}
    \small Рисунок 5 --- Предсказания мелкой сети в PCA пространстве
\end{figure}

\subsection*{Матрицы ошибок и метрики}

\begin{lstlisting}
from sklearn.metrics import confusion_matrix
import seaborn as sns

# Матрица ошибок для глубокой сети
cm = confusion_matrix(y_test, y_pred_deep_cls)

fig, ax = plt.subplots(figsize=(8, 6))
sns.heatmap(cm, annot=True, fmt='d', cmap='Blues', ax=ax, 
            xticklabels=['Class 0', 'Class 1', 'Class 2'],
            yticklabels=['Class 0', 'Class 1', 'Class 2'])
ax.set_ylabel('True Label')
ax.set_xlabel('Predicted Label')
ax.set_title('Матрица ошибок Deep Network')
plt.tight_layout()
plt.savefig('fig_6.png', dpi=150, bbox_inches='tight')
plt.show()
\end{lstlisting}

\begin{figure}[h!]
    \centering
    \includegraphics[width=0.85\textwidth]{fig_6.png}
    \small Рисунок 6 --- Матрица ошибок для глубокой нейронной сети
\end{figure}

\subsection*{Распределение вероятностей предсказаний}

\begin{lstlisting}
# Анализ вероятностей предсказаний
max_probs_shallow = np.max(y_pred_shallow, axis=1)
max_probs_medium = np.max(y_pred_medium, axis=1)
max_probs_deep = np.max(y_pred_deep, axis=1)

fig, ax = plt.subplots(figsize=(10, 6))
ax.hist(max_probs_shallow, bins=20, alpha=0.5, label='Shallow', color='#FF6B6B')
ax.hist(max_probs_medium, bins=20, alpha=0.5, label='Medium', color='#4ECDC4')
ax.hist(max_probs_deep, bins=20, alpha=0.5, label='Deep', color='#45B7D1')
ax.set_xlabel('Максимальная вероятность')
ax.set_ylabel('Количество образцов')
ax.set_title('Распределение максимальных вероятностей предсказаний')
ax.legend()
plt.tight_layout()
plt.savefig('fig_7.png', dpi=150, bbox_inches='tight')
plt.show()
\end{lstlisting}

\begin{figure}[h!]
    \centering
    \includegraphics[width=0.85\textwidth]{fig_7.png}
    \small Рисунок 7 --- Распределение максимальных вероятностей предсказаний
\end{figure}

\section*{Оценка качества}

\section*{Результаты}

\begin{center}
\begin{tabular}{|l|c|}
\hline
\textbf{Архитектура} & \textbf{Test Accuracy} \\
\hline
Мелкая (2 слоя) & 0.9777 \\
\hline
Средняя (3 слоя) & 0.9889 \\
\hline
Глубокая (5 слоев) & 1.0000 \\
\hline
\end{tabular}
\end{center}

\section*{Выводы}

\begin{enumerate}
    \item Глубокие сети обеспечивают более точное предсказание благодаря большей емкости
    \item Dropout-слои критичны для предотвращения переобучения
    \item Для датасета Iris даже мелкая сеть показывает хорошие результаты
    \item Глубокие сети требуют больше данных для избежания переобучения
    \item Батч-нормализация может улучшить сходимость обучения
    \item t-SNE лучше сохраняет локальную структуру данных чем PCA
    \item Максимальная вероятность предсказания может использоваться как мера уверенности модели
\end{enumerate}

\end{document}
